\documentclass[10pt,a4paper]{amsart}
\usepackage{amsmath,amssymb,amsthm,graphicx}
% Retained minimal preamble for the Stacks Project chapter varieties.tex.
% Verbatim excerpt of the chapter preamble: Xy-pic setup (native preamble lines 14-18), the native theorem
% declarations with the mathematical macros (native preamble lines 154-217), and the project own no-op
% definitions of the reference/slogan/history display environments (native wrapper.tex lines 1-3) that the
% counted statement uses for its slogan. Copied mechanically; nothing added, removed or redefined.
% For commutative diagrams we use Xy-pic
\usepackage[all]{xy}

% We use 2cell for 2-commutative diagrams.
\xyoption{2cell}
% Theorem environments.
%
\theoremstyle{plain}
\newtheorem{theorem}[subsection]{Theorem}
\newtheorem{proposition}[subsection]{Proposition}
\newtheorem{lemma}[subsection]{Lemma}

\theoremstyle{definition}
\newtheorem{definition}[subsection]{Definition}
\newtheorem{example}[subsection]{Example}
\newtheorem{exercise}[subsection]{Exercise}
\newtheorem{situation}[subsection]{Situation}

\theoremstyle{remark}
\newtheorem{remark}[subsection]{Remark}
\newtheorem{remarks}[subsection]{Remarks}

\numberwithin{equation}{subsection}

% Macros
%
\def\lim{\mathop{\mathrm{lim}}\nolimits}
\def\colim{\mathop{\mathrm{colim}}\nolimits}
\def\Spec{\mathop{\mathrm{Spec}}}
\def\Hom{\mathop{\mathrm{Hom}}\nolimits}
\def\Ext{\mathop{\mathrm{Ext}}\nolimits}
\def\SheafHom{\mathop{\mathcal{H}\!\mathit{om}}\nolimits}
\def\SheafExt{\mathop{\mathcal{E}\!\mathit{xt}}\nolimits}
\def\Sch{\mathit{Sch}}
\def\Mor{\mathop{\mathrm{Mor}}\nolimits}
\def\Ob{\mathop{\mathrm{Ob}}\nolimits}
\def\Sh{\mathop{\mathit{Sh}}\nolimits}
\def\NL{\mathop{N\!L}\nolimits}
\def\CH{\mathop{\mathrm{CH}}\nolimits}
\def\proetale{{pro\text{-}\acute{e}tale}}
\def\etale{{\acute{e}tale}}
\def\QCoh{\mathit{QCoh}}
\def\Ker{\mathop{\mathrm{Ker}}}
\def\Im{\mathop{\mathrm{Im}}}
\def\Coker{\mathop{\mathrm{Coker}}}
\def\Coim{\mathop{\mathrm{Coim}}}

% Boxtimes
%
\DeclareMathSymbol{\boxtimes}{\mathbin}{AMSa}{"02}

%
% Macros for moduli stacks/spaces
%
\def\QCohstack{\mathcal{QC}\!\mathit{oh}}
\def\Cohstack{\mathcal{C}\!\mathit{oh}}
\def\Spacesstack{\mathcal{S}\!\mathit{paces}}
\def\Quotfunctor{\mathrm{Quot}}
\def\Hilbfunctor{\mathrm{Hilb}}
\def\Curvesstack{\mathcal{C}\!\mathit{urves}}
\def\Polarizedstack{\mathcal{P}\!\mathit{olarized}}
\def\Complexesstack{\mathcal{C}\!\mathit{omplexes}}
% \Pic is the operator that assigns to X its picard group, usage \Pic(X)
% \Picardstack_{X/B} denotes the Picard stack of X over B
% \Picardfunctor_{X/B} denotes the Picard functor of X over B
\def\Pic{\mathop{\mathrm{Pic}}\nolimits}
\def\Picardstack{\mathcal{P}\!\mathit{ic}}
\def\Picardfunctor{\mathrm{Pic}}
\def\Deformationcategory{\mathcal{D}\!\mathit{ef}}
\newenvironment{reference}{}{}
\newenvironment{slogan}{}{}
\newenvironment{history}{}{}
\usepackage{cleveref}
\crefname{originalreference}{original reference}{original references}
\Crefname{originalreference}{Original reference}{Original references}
\makeatletter
\newlabel{algebra-lemma-integral-closure-Dedekind}{{09IG}{source}{}{sourcepointer-1}{}}
\newlabel{algebra-lemma-integral-closure-Dedekind@cref}{{[originalreference][0][]09IG}{[1][0][]source}{}{}{}}
\makeatother
\crefname{originalreference}{original reference}{original references}
\Crefname{originalreference}{Original reference}{Original references}
\makeatletter
\newlabel{algebra-lemma-characterize-finite-in-terms-of-integral}{{02JJ}{source}{}{sourcepointer-2}{}}
\newlabel{algebra-lemma-characterize-finite-in-terms-of-integral@cref}{{[originalreference][0][]02JJ}{[1][0][]source}{}{}{}}
\makeatother
\crefname{originalreference}{original reference}{original references}
\Crefname{originalreference}{Original reference}{Original references}
\makeatletter
\newlabel{algebra-lemma-integral-overring-surjective}{{00GQ}{source}{}{sourcepointer-3}{}}
\newlabel{algebra-lemma-integral-overring-surjective@cref}{{[originalreference][0][]00GQ}{[1][0][]source}{}{}{}}
\makeatother
\title{Powers of an element of the Jacobson radical of the integral closure of a one-dimensional semi-local domain land back in the domain}
\author{The Stacks Project Authors}
\date{stacks-project:varieties.tex:local-snapshot-20261011; GNU FDL 1.2 or later}
\begin{document}
\maketitle
\noindent\textit{Editorial note.} All delivered mathematical text is the author's own, copied verbatim from the pinned Stacks Project chapter file varieties.tex; no statement or proof step was rewritten, reordered or removed.\par\medskip
\noindent\textit{Editorial note.} The counted span cites three cross-chapter labels, each recorded in assembly.external\_references with its canonical tag as a navigation pointer.\par\medskip
\noindent\textit{Editorial note.} The proof cites no bibliography item inside the counted span.\par\medskip
\noindent\textit{Allowed standard prerequisites (editorial).}\begin{itemize}
\item Structure of the integral closure of a one-dimensional Noetherian domain. Stacks Project, Algebra, Lemma algebra-lemma-integral-closure-Dedekind (tag 034W). Gives that B' is a semi-local Dedekind domain.
\item Finiteness criteria for integral extensions and surjectivity of the spectrum map. Stacks Project, Algebra, Lemmas algebra-lemma-characterize-finite-in-terms-of-integral (tag 00GM) and algebra-lemma-integral-overring-surjective (tag 00GT). Make the intermediate ring A finite over B and show that A is semi-local.
\end{itemize}\medskip
\ifdefined\hypertarget\hypertarget{sourcepointer-1}{}\fi
\noindent\textit{Source reference.} Original reference 09IG: sources/stacks-chow/tags.txt, source lines 12039--12039. Cross-chapter standard prerequisite cited by the author inside the counted span; it is used as a pointer and its own proof is not part of this item.\par\medskip
\ifdefined\hypertarget\hypertarget{sourcepointer-2}{}\fi
\noindent\textit{Source reference.} Original reference 02JJ: sources/stacks-chow/tags.txt, source lines 3338--3338. Cross-chapter standard prerequisite cited by the author inside the counted span; it is used as a pointer and its own proof is not part of this item.\par\medskip
\ifdefined\hypertarget\hypertarget{sourcepointer-3}{}\fi
\noindent\textit{Source reference.} Original reference 00GQ: sources/stacks-chow/tags.txt, source lines 616--616. Cross-chapter standard prerequisite cited by the author inside the counted span; it is used as a pointer and its own proof is not part of this item.\par\medskip
\begin{lemma}
\label{lemma-semi-local-dimension-one-conductor}
Let $B$ be a semi-local Noetherian domain of dimension $1$.
Let $B'$ be the integral closure of $B$ in its fraction field.
Then $B'$ is a semi-local Dedekind domain.
Let $x$ be a nonzero element of the Jacobson radical of $B'$.
Then for every $y \in B'$ there exists an $n$ such that
$x^n y \in B$.
\end{lemma}

\begin{proof}
Let $\mathfrak m_B$ be the Jacobson radical of $B$.
The structure of $B'$ results from
Algebra, Lemma \ref{algebra-lemma-integral-closure-Dedekind}.
Given $x, y \in B'$ as in the statement of the lemma consider
the subring $B \subset A \subset B'$ generated by $x$ and $y$.
Then $A$ is finite over $B$ (Algebra, Lemma
\ref{algebra-lemma-characterize-finite-in-terms-of-integral}).
Since the fraction fields of $B$ and $A$ are the same we see
that the finite module $A/B$ is supported on the set of
closed points of $B$. Thus $\mathfrak m_B^n A \subset B$ for
a suitable $n$. Moreover, $\Spec(B') \to \Spec(A)$ is
surjective (Algebra, Lemma \ref{algebra-lemma-integral-overring-surjective}),
hence $A$ is semi-local as well. It also follows that
$x$ is in the Jacobson radical $\mathfrak m_A$ of $A$. Note that
$\mathfrak m_A = \sqrt{\mathfrak m_B A}$.
Thus $x^m y \in \mathfrak m_B A$ for some $m$.
Then $x^{nm} y \in B$.
\end{proof}

\end{document}
